\hypertarget{Introduction}{%
\chapter{Introduction}\label{Introduction}}

\section{Présentation de la structure
d'accueil}

Durant la période de mon stage , j'ai été accueilli au
\textbf{Laboratoire de Mathématiques Informatique et Application
(LAMIA)} de l'Université des Antilles (UA).

Pour présenter cette structure, il me faut tout d'abord présenter
l'université à laquelle il est rattaché.

\hypertarget{luniversite-des-antilles}{%
\subsection{L'université des Antilles}\label{luniversite-des-antilles}}

Bien que ce soit l'université dans laquelle j'ai fait toutes mes études,
voici quelques chiffres que je ne connaissais pas et qui donnent la
mesure de sa taille :

L'Université des Antilles s'organise autour deux pôles universitaires
régionaux autonomes : le « Pôle Guadeloupe » et le « Pôle Martinique ».

Sur ces pôles, l'Université assure des missions d'\emph{enseignement} et
de \emph{recherche}, assistées par des \emph{administratifs et des
techniciens}.

\hypertarget{administration-et-personnel-technique}{%
\subsubsection{Administration et personnel
technique}
\label{administration-et-personnel-technique}}

l'UA emploie 414 Administratifs et Techniciens (environ 200 personnes
pour l'administation centrale et 100 répartis sur chaque pôle)

\hypertarget{enseignements}{%
\subsubsection{Enseignements}\label{enseignements}}

L'UA délivre des diplomes de la licence au doctorat dans de nombreux
domaines. Au total, cela représente :

\begin{itemize}
\tightlist
\item
  484 enseignants-chercheurs (environ 240 pour chaque pôle)
\item
  12 000 étudiants (environ 7000 pour la Guadeloupe , 5000 pour la
  Martinique)
\end{itemize}

Pour l'informatique, cela représente : - autour de 20
enseignants-chercheurs - autour de 120 étudiants

\hypertarget{le-lamia}{%
\subsection{Le LAMIA}\label{le-lamia}}

Le \textbf{Laboratoire de Mathématiques Informatique et Application
(LAMIA)}, comme son nom l'indique, se concentre sur les recherches en
informatiques et mathématiques.

Il compte une soixantaine de membres (Professeurs des Universités,
Maitres de Conférences, ATER, Doctorants) répartis sur deux pôles
(Guadeloupe et Martinique) au sein de trois équipes internes :

\begin{itemize}
\tightlist
\item
  Equipe
  \href{http://lamia.univ-ag.fr/index.php?page=equipe-mathematiques}{\textbf{Mathématiques}
  (analyse variationnelle, analyse numérique, EDP, analyse statistique,
  mathématiques discrètes)} ;
\item
  Equipe Informatique
  \href{http://lamia.univ-ag.fr/index.php?page=equipe-danais}{\textbf{DANAIS}
  : Data analytics and big data gathering with sensors} ;
\item
  Equipe Informatique
  \href{http://lamia.univ-ag.fr/index.php?page=equipe-aid}{\textbf{AID}
  : Apprentissages Interactions Donnees} ;
\end{itemize}

De plus, le LAMIA accueille en son sein un groupe de chercheurs associés
travaillant en Epidémiologie clinique et médecine.



\hypertarget{contexte-guxe9nuxe9ral}{%
\section{Contexte général}\label{contexte-guxe9nuxe9ral}}


\hypertarget{Etat de l'art}{%
\section{Etat de l'art}\label{Etat de l'art}


\hypertarget{Méthodologie}{
\section{Méthodologie}\label{Méthodologie}}
\subsection{Outils utilisés}



\subsubsection{Présentation de Python}

\begin{figure}[h]
  \begin{center}
  \includegraphics[width=4cm]{./images/Python_Logo.png}\label{Python}
  \caption{Logo de Python.}
  \end{center}
\end{figure}

Python est un langage de programmation interprété\footnote{\label{interprete}Langage nécéssitant un programme informatique qui joue le rôle d’interface entre le projet et le processeur appellé interpréteur, pour exécuter du code.} qui sera utiliser pour l'ensemble du projet.

\subsubsection{Présentation Qt}

\begin{figure}[h]
  \begin{center}
  \includegraphics[width=2cm]{./images/Qt_logo_2016.png}\label{Qt}
  \caption{Logo de Qt.}
  \end{center}
\end{figure}

QT est une librairie\footnote{\label{librairie}Une librairie est un fichier contenant du code (généralement un ensemble de fonction et classes permettant de faciliter et/où de réaliser certain programme)} qui permet le création d'interface graphique en Python.Que nous utiliserons pour créer l'interface que l'on utilisera au cours du projet.

\subsubsection{Présentation Cplex}

\begin{figure}[h]
  \begin{center}
  \includegraphics[width=4cm]{./images/cplex.png}\label{Cplex}
  \caption{Logo de Cplex.}
  \end{center}
\end{figure}

Cplex est une librairie\footref{librairie} qui permet la modélisation et la résolution de problème d'optimisation linéaire\footnote{Terme que j'expliquerais plus tard dans mon rapport}.
\newline
\newline
\subsubsection{Présentation de GitHub}
\begin{figure}[h]
  \begin{center}
  \includegraphics[width=3cm]{./images/github.jpg}\label{GitHub}
  \caption{Logo de GitHub.}
  \end{center}
\end{figure}


Nous pouvons définir GitHub comme une plateforme de développement de projet informatique en groupe. Elle simplifie grandement le développement de projets. Elle permet de versioner ses programmes et d'y apporter des modifications en temps réel à plusieurs.

\subsubsection{Présentation de LaTex}

\begin{figure}[h]
  \begin{center}
\includegraphics[width=2cm]{./images/Latex.png}\label{LaTex}
\caption{Logo de LaTex}
\end{center}
\end{figure}

Nous pouvons dire que LaTex est un langage de traitement de texte tel que le markdown qui permet de mettre en forme notre texte de manière \"scientifique\" cela veut dire que. LaTex permet une faciliter d'écriture des équations et de toutes les écriture mathématiques.Permet de par ses nombreux package une quasi-infinité de possibilitées.
L'utilisation de cet outil permettra une synergie entre ceux-ci car LaTex peut-être utiliser avec un simple bloc-note c'est donc du texte ce qui permet une intéraction facilitée avec GitHub d'ailleurs ce rapport est écrit avec Latex et retrouvable sur GitHub.\newline


\hypertarget{Annonce du plan}{%
\section{Annonce du plan}\label{annonce du plan}}

\subsection{Présentations des stratégies de résolution}

Dans cette première partie je commencerais par vous présenter la première solution de résolution choisie qui est la résolution du sudoku en tant que problème d'optimisation linéaire.
Je vous expliquerai ce qu'est un problème d'optimisation linéaire puis vous décrirai l'algorithme du cplex qui nous permettra de le résoudre.
En deuxième grande sous partie de cette section je vous présenterai l'algorithme du backtraking.
Je détaillerai son fonctionnement puis vous expliquerais en quoi c'est un algorithme que nous pouvons constament améliorer.
Et finalement je terminerais par expliquer pourquoi avoir choisi ces deux méthodes de résolutions.
\subsection{Implémentation des stratégies}
La première sous partie de cette grande sous partie commencera avec la modélisation et l'implémentation d'un sudoku en python et l'implémentation de l'interface graphique.

\subsection{Test du résolveur}
